\documentclass[10pt,a4paper]{amsart}
\usepackage[margin=19mm]{geometry}
\usepackage{amsmath,amssymb,mathtools}
\usepackage[scr=rsfs,cal=euler]{mathalfa}
\usepackage{xfrac}
\usepackage[unicode,hidelinks,bookmarks=false]{hyperref}
\newcommand{\ie}{i.e.\ }
\newcommand{\cf}{cf.\ }
\newcommand{\resp}{respectively\ }
\newcommand{\Subsection}{\subsection}
\providecommand{\nobreakdash}{\nobreak\hbox{-}}
\setlength{\emergencystretch}{2em}
\multlinegap0pt
\usepackage{amssymb}
\usepackage{enumerate}
\usepackage{tikz}
\newtheorem{lem}{Lemma}
\title{Licci Property of Complementary Edge Ideals}
\author{Vivek Bhabani Lama}
\date{Source: arXiv:2603.16193v1 (CC BY 4.0)}
\begin{document}
\maketitle

\noindent\emph{Source and licence.} Licci Property of Complementary Edge Ideals. Version arXiv:2603.16193v1 (CC BY 4.0), distributed by arXiv under the Creative Commons Attribution 4.0 International licence, http://creativecommons.org/licenses/by/4.0/, as recorded in the arXiv licence metadata for this exact version and shown on the abstract page https://arxiv.org/abs/2603.16193v1. Full licence notice: \texttt{licenses/arxiv-2603.16193-CC-BY-4.0.txt}.\par\smallskip

\noindent\textit{Editorial scope.} Counted unit: the article's lem height (source0.tex lines 193--195). The article's complete original proof of it is delivered with it (source0.tex lines 196--198, one \texttt{proof} environment of 3 lines). No mathematics is written, repaired or re-derived: the statement and the proof are the article's own lines, in the article's own notation, and no retained line is substituted or re-flowed.  No further source-local context is needed: every cross-reference of every retained line resolves inside the two delivered spans.  Nothing was omitted from any retained span, so the package carries no omission record.  No retained line contains a citation, so the wrapper carries no bibliography and no bibliography entry.  No retained line contains a figure environment, an image-inclusion command or drawing code, so no image is delivered and the package declares no asset dependency.  Editorial formatting only, in addition: an AMS \texttt{amsart} preamble that restates the article's own declarations and macros used by the retained lines, namely the environments \texttt{lem} on separate counters, together with the standard packages those lines need.  The retained environments print in this document as Lemma 1 in order of appearance, whereas the article prints the same items with the numbering of its own full document.  The complete native source of the article as distributed in the arXiv e-print for this exact version is bundled unchanged as \texttt{sources/196526/source0.tex}.
\begin{lem}\label{lemma on height}
Let $n\geq 3$ be a positive integer and $G=K_n$ be a complete graph on the vertex set $\{1,2,\dots,n\}$. Then, $\mathrm{ht}(I_c(G))=3$.   
\end{lem}

\begin{proof}
 We observe that for a complete graph $G=K_n$, the ideal $I_c(G)$ is an ideal generated by all possible squarefree monomials of degree $n-2$. Also, for a monomial ideal $J$ generated by all possible squarefree monomials of a fixed degree $d$, it is well-known that the height of $J$ is given by $n-d+1$. Therefore, it easily follows that $\mathrm{ht}(I_c(G))=n-(n-2)+1=3$.     
\end{proof}

\end{document}
